\section{Conclusioni}

\subsection{Milestone 1}

La Milestone 1 ha prodotto un dataset di 15.338 snapshot di classi Java per Apache OpenJPA,
corredato di 19 metriche e dell'etichetta di bugginess. Le scelte metodologiche (filtraggio
versioni, sanitizzazione ticket, gestione versioni senza tag Git) privilegiano
la coerenza dei dati rispetto alla massimizzazione della dimensione del dataset.

\subsection{Milestone 2}

\begin{description}
  \item[\textit{Quale classificatore è più accurato degli altri?}]
  RF (RandomForest + No FS + SMOTE) domina su tutte e cinque le metriche nella sua
  configurazione ottimale: Kappa $0{,}61$, AUC $0{,}94$, NPofB20 $0{,}88$, Precision
  $0{,}63$, Recall $0{,}66$ (Tabelle~\ref{tab:r2-rf-full}--\ref{tab:r2-ibk-full}).
  NaiveBayes non supera Kappa $0{,}29$ per violazione dell'indipendenza condizionale;
  IBk si ferma a Kappa $0{,}47$ e NPofB20 $0{,}64$.

  \item[\textit{Il miglior classificatore cambia in base a dataset / numero di release / metrica?}]
  \textbf{Per dataset}: è stato analizzato un solo progetto (Apache OpenJPA); il confronto
  tra dataset richiederebbe l'abilitazione di ulteriori progetti in \texttt{config.yaml},
  che la struttura del CSV di output supporta senza modifiche al codice.
  \textbf{Per numero di release}: questo test non è stato effettuato --- il progetto usa
  10$\times$10-fold CV sull'intero dataset con le release mescolate casualmente tra i fold.
  Per approssimare l'effetto del numero di release in training basterebbe variare il numero
  di fold (es.\ 5-fold usa l'80\% dei dati in training, 2-fold il 50\%); un controllo
  rigoroso per release richiederebbe invece una strategia Walk-Forward (train su release
  1, test su 2; poi train su 1--2, test su 3; ecc.).
  \textbf{Per metrica}: RF è il migliore su ogni singola metrica --- AUC $0{,}94$ vs
  $0{,}80$, Kappa $0{,}61$ vs $0{,}47$, NPofB20 $0{,}88$ vs $0{,}64$ rispetto a NB e IBk.
  La scelta del classificatore non cambia al variare della metrica di riferimento.

  \item[\textit{Quale è il classificatore migliore per una determinata metrica?}]
  RF è il migliore su AUC, Kappa, NPofB20 e Precision.
  Per il solo Recall, RF + Undersampling raggiunge $0{,}87$ ma Kappa crolla a $0{,}37$ e
  Precision a $0{,}31$: il guadagno è illusorio su un dataset al 9\% di buggy.
  In nessuna configurazione un classificatore diverso da RF supera RF su Kappa o NPofB20
  (Tabelle~\ref{tab:r2-rf-full}--\ref{tab:r2-ibk-full}).
\end{description}

\subsection{Milestone 3}

\begin{description}
  \item[\textit{Is BClassifier accurate?}]
  Sì. In 10$\times$10-fold cross-validation su A, RF + No FS + SMOTE ottiene Precision
  $0{,}63$, Recall $0{,}66$, AUC $0{,}94$, Kappa $0{,}61$
  (Tabella~\ref{tab:r2-rf-full}).

  \item[\textit{How many buggy classes could have been prevented by having zero smells?}]
  I risultati sono in Tabella~\ref{tab:r3-prevention}.
  \begin{itemize}[noitemsep]
    \item \textbf{In total}: \textbf{63} classi smettono di essere predette buggy azzerando NSmells.
    \item \textbf{In proportion}: il \textbf{4,57\%} delle 1.378 classi effettivamente buggy
          nel dataset A.
    \item \textbf{Out of the preventable ones}: il \textbf{5,60\%} delle 1.125 classi predette
          buggy in B+ (quelle con almeno uno smell che il modello giudica difettose).
  \end{itemize}
  Il dato più significativo è l'effetto inverso: azzerare NSmells genera anche 75 istanze
  che passano da CLEAN a BUGGY, per un bilancio netto di $+12$ bug predetti in B rispetto
  a B+ (Tabella~\ref{tab:r3-datasets}). La causa è la forte correlazione tra NSmells e le
  metriche di processo ($\rho = 0{,}78$ con LOC, Tabella~\ref{tab:r3-correlation}): azzerare
  NSmells distorce le distribuzioni delle feature correlate, spingendo il modello verso
  predizioni inattese.
\end{description}

\subsection{Milestone 4}

\begin{description}
  \item[\textit{Does C\textsubscript{X} compile?}]
  \textbf{[TODO]} Riportare se $C_1$--$C_4$ compilano da soli e integrati con il sistema.

  \item[\textit{Does C\textsubscript{X} have smells?}]
  \textbf{[TODO]} Indicare per ogni versione: nessuno smell residuo, smell vecchi mantenuti,
  o nuovi smell introdotti dal refactoring.

  \item[\textit{Is any feature positively correlated with bugginess higher in C\textsubscript{X} than in C\textsubscript{0}?}]
  \textbf{[TODO]} Se sì, il refactoring potrebbe non aver migliorato la manutenibilità.

  \item[\textit{Is any feature negatively correlated with bugginess higher in C\textsubscript{X} than in C\textsubscript{0}?}]
  \textbf{[TODO]} Se sì, il refactoring ha probabilmente migliorato la manutenibilità.
\end{description}

